% Full HSP constructions, without changing the existing equational proofs.

\begin{annexitem}{annex:p10-hsp-preservation}{thm:p10-hsp-preservation}{Preservation and equality of identity theories}
\begin{proof}
Use the existing term-evaluation induction in
\cref{thm:p10-term-universal}.
If \(\mathbf B\leq\mathbf A\), evaluation of a term
at arguments from \(B\) agrees with its evaluation in
\(\mathbf A\), because the restricted operations agree.
This includes constants, which belong to every subalgebra.
Thus every identity of \(\mathbf A\) holds in \(\mathbf B\).

Let \(h:\mathbf A\to\mathbf B\) be onto, and suppose
\(\mathbf A\models p\approx q\).
Only finitely many variables occur in \(p,q\).
For arbitrary values of these variables in \(B\), select
preimages in \(A\); finite choice suffices.
Evaluation commutes with \(h\), so the equal values of
\(p,q\) in \(\mathbf A\) have equal images in
\(\mathbf B\). This proves the identity for every
valuation in \(\mathbf B\). A ground identity needs
no variable preimages; constant preservation is enough.

For a set-indexed product, induction on terms shows
that term evaluation is coordinatewise. If every factor
satisfies \(p\approx q\), the two values have equal
coordinates everywhere and are equal.
Under ZFC a product of nonempty factors has a nonempty
carrier. For the empty product there is one empty tuple,
so every pair of term values is equal, including ground
terms. Identities also survive isomorphisms by applying
the onto-homomorphism argument.

Repeated preservation proves that every identity of \(K\)
holds in \(\mathsf H\mathsf S\mathsf P(K)\).
Conversely every \(\mathbf A\in K\) is an onto
image of a subalgebra of its one-factor product, so
\(K\subseteq\mathsf H\mathsf S\mathsf P(K)\).
An identity of that larger class is consequently an
identity of \(K\). These two inclusions prove (HSP.2),
including \(K=\varnothing\).
For nonempty \(K\), the same one-factor inclusion and
preservation argument apply to the book's \(HSP\).
\end{proof}
\end{annexitem}

\begin{annexitem}{annex:p10-hsp-witness-product}{thm:p10-hsp-witness-product}{The kernel set, separating product and relative-free map}
\begin{proof}
Put \(\mathbf T=\mathbf T(X)\).
Every homomorphism \(h:\mathbf T\to\mathbf A\)
has a congruence kernel. Its image is a nonempty
subalgebra: a term exists, operations are preserved,
and constants are preserved.
The map
\[
 \mathbf T/\ker h\longrightarrow h(T),\qquad
 [p]_{\ker h}\longmapsto h(p)
\]
is well defined, injective, onto and operation preserving.
Its inverse preserves each operation by taking preimages
of a finite argument tuple. It is therefore an isomorphism.
If \(\mathbf A\in K\), this gives
\(\ker h\in\Phi_K(X)\).
Conversely if \(\phi\in\Phi_K(X)\), the definition
supplies an embedding \(j_\phi\) of
\(\mathbf T/\phi\) into some \(\mathbf A_\phi\in K\).
Then \(h_\phi=j_\phi q_\phi\) has kernel exactly \(\phi\).
Hence \(\Phi_K(X)\) is exactly the set of kernels
obtained from homomorphisms into members of \(K\).
This proves (HSP.4) without forming a product over a
proper class of all such homomorphisms.

The congruence set \(\operatorname{Con}\mathbf T\)
is a subset of \(\mathcal P(T^2)\); separation gives
the set \(\Phi_K(X)\). Its intersection is a congruence.
To see it is nonempty as an indexing set, choose
\(\mathbf A_0\in K\) and one \(a_0\in A_0\).
The constant map \(X\to A_0\) extends to term evaluation,
also when \(X=\varnothing\) and ground terms exist.
Its kernel belongs to \(\Phi_K(X)\).

If \(K\) is proper, its membership condition is the
declared class formula. For each \(\phi\) in the set
\(\Phi_K(X)\), there exists a pair
\((\mathbf A,j)\) of the required kind.
Collection supplies a set of such witnesses covering
every \(\phi\). Tag each witness with its \(\phi\);
choice on the resulting nonempty set fibers selects
one \((\mathbf A_\phi,j_\phi)\) per \(\phi\).
Thus the product in (HSP.5) is set-indexed.
No choice over the whole proper class \(K\) is made.

For the map \(d\), if \([p]_\theta=[q]_\theta\),
then \((p,q)\) lies in every \(\phi\), whence
\(h_\phi(p)=h_\phi(q)\) for every coordinate.
This proves well-definedness. Each coordinate is a
homomorphism, so quotient and product operation clauses
show that \(d\) preserves every basic operation,
including nullary ones.
If \(d([p])=d([q])\), each coordinate equality says
\((p,q)\in\ker h_\phi=\phi\).
Their intersection is \(\theta\); thus \([p]=[q]\).
So \(d\) is injective. Its image \(\mathbf B\)
is a nonempty subalgebra of the product, and the
inverse \(d^{-1}:\mathbf B\to\mathbf F_K(\bar X)\)
preserves operations by the same image argument above.
This proves the asserted \(ISP(K)\) membership.
In fact a fixed term \(p_0\in T\) provides the tuple
\((h_\phi(p_0))_\phi\), independently exhibiting a
point of this particular product.

The term algebra is generated by \(X\), with ground
terms included through the constants. Its quotient
and the image of \(d\) are therefore generated by
the indicated generator images, by induction on terms.
Given \(v:X\to A\), let \(e_v:\mathbf T\to\mathbf A\)
be term evaluation. Its kernel belongs to \(\Phi_K(X)\).
Hence \(\theta\subseteq\ker e_v\), so
\(\bar v([p])=e_v(p)\) is well defined and preserves
operations. It extends \(v\).
Every element of the quotient is a term class,
so induction on terms gives uniqueness.

Finally, every member of \(K\) satisfies each common
identity, whatever its variable alphabet. Products,
subalgebras and isomorphic copies preserve that
identity by \cref{thm:p10-hsp-preservation}.
Since \(\mathbf F_K(\bar X)\) is explicitly an
isomorphic copy of \(\mathbf B\) in the product,
it satisfies each such identity under every valuation.
This conclusion is not restricted to the canonical
values \([x]\).
\end{proof}
\end{annexitem}

\begin{annexitem}{annex:p10-hsp-representation}{thm:p10-hsp-representation}{The target quotient and both directions of HSP}
\begin{proof}
First let \(K\ne\varnothing\) and
\(\mathbf A\models\operatorname{Id}_K(W)\).
Take the tagged \(X=\{x_a:a\in A\}\).
It is nonempty; therefore \(\mathbf T(X)\) exists
even without constants. Term evaluation \(e_A\)
extends \(x_a\mapsto a\) and is onto.

We prove \(\theta_K(X)\subseteq\ker e_A\).
Suppose \((p,q)\in\theta_K(X)\), and let
\(Y\subseteq X\) be the finite union of their variable
supports. Choose an injective renaming
\(\rho:Y\to W\). Replace those variables to form
the terms \(p_\rho,q_\rho\) over \(W\).
For arbitrary \(\mathbf C\in K\) and
valuation \(u:W\to C\), choose one \(c_0\in C\).
Set \(v(y)=u(\rho(y))\) for \(y\in Y\) and
\(v(x)=c_0\) on the remaining \(X\).
The evaluation homomorphism \(e_v\) has
\((p,q)\) in its kernel by (HSP.4).
The term-substitution law therefore gives
\[
 p_\rho^{\mathbf C}(u)
    =p^{\mathbf C}(v)
    =q^{\mathbf C}(v)
    =q_\rho^{\mathbf C}(u).
 \]
Thus \(p_\rho\approx q_\rho\in\operatorname{Id}_K(W)\).
This argument also applies to \(Y=\varnothing\):
the terms are ground and no renaming values are needed.

Choose one \(a_0\in A\), and give the distinct
variables \(\rho(x_a)\), for \(x_a\in Y\), the
values \(a\), using \(a_0\) for all other variables
of \(W\). The identity in \(\mathbf A\)
and the same substitution law imply
\(e_A(p)=e_A(q)\), as required.
The finite-support argument is why a countably
infinite \(W\) suffices even when \(X\) or \(\mathcal F\)
is uncountable.

Consequently \(r_A([p])=e_A(p)\) is well defined.
The quotient operations and homomorphism property
of \(e_A\) prove that it is a homomorphism, including
on constants. For every \(a\in A\),
\(r_A([x_a])=a\); it is onto.
Use the embedding \(d\) and image \(\mathbf B\)
of \cref{thm:p10-hsp-witness-product}.
The map \(b_A=r_A d^{-1}\) is an onto homomorphism
from that concrete subalgebra to \(\mathbf A\).
Equivalently the map
\[
 \mathbf B/\ker b_A\longrightarrow\mathbf A,\qquad
 [b]\longmapsto b_A(b)
 \]
is well defined, injective, onto and a homomorphism,
with an operation-preserving inverse.
This supplies the quotient form of the representation.
Therefore \(\mathbf A\in HSP(K)\), using a
nonempty indexing set, and also in the host's
\(\mathsf H\mathsf S\mathsf P(K)\).
The opposite containment in \(M(\operatorname{Id}_K(W))\)
follows from preservation. This proves (HSP.8)
for nonempty \(K\) and the equivalence of the two
product conventions there.

For \(K=\varnothing\), every identity over \(W\)
is in \(\operatorname{Id}_K(W)\). In particular
take \(x\approx y\) for distinct variables.
For any \(a,b\) in a model, a valuation sending
them to \(x,y\) yields \(a=b\); the carrier is
one-element since it is nonempty.
Conversely every one-element algebra satisfies
all identities. The host's empty product produces
the one-element algebra; its nonempty subalgebras
and homomorphic images are one-element, and every
isomorphic one-element copy is a homomorphic image.
Thus its HSP closure of \(\varnothing\) is
exactly \(M(\operatorname{Id}_\varnothing(W))\).
With the book's nonempty products,
\(HSP(\varnothing)=\varnothing\); this case cannot
be included in that equality.

Finally let a class \(\mathcal C=M(\Sigma)\)
be equational. Preservation gives all the
host closure conditions, and its one-element
algebra supplies nonemptiness.
Conversely if \(\mathcal C\) is closed under
the host H/S/P, it contains the empty product.
The three successive closure conditions give
\(\mathsf H\mathsf S\mathsf P(\mathcal C)\subseteq\mathcal C\).
The one-factor inclusion gives the reverse.
Apply (HSP.8) with \(K=\mathcal C\):
\[
 \mathcal C=M(\operatorname{Id}_{\mathcal C}(W)).
 \]
This is an equational axiomatization by a set.
Under the book convention, assume instead
\(\mathcal C\ne\varnothing\) and H/S/P closure.
Any member has an onto homomorphism to the
one-element algebra, so that algebra already
belongs to \(\mathcal C\). Empty-product closure
then follows, and the same proof applies.
\end{proof}
\end{annexitem}

\begin{annexitem}{annex:p10-hsp-original-product}{thm:p10-hsp-original-product}{The separately declared 1935 all-valuations route}
\begin{proof}
Since \(\mathfrak A\) and \(X\) are sets, the
index collection of pairs \((i,\xi)\), with
\(\xi:X\to A_i\), is a set: the union of the
carriers is a set, and all these maps can be
bounded by a power set of \(X\) times that union.
For each \(\xi\), the universal property of terms
gives \(\widehat\xi:\mathbf T(X)\to\mathbf A_i\).
Its image is a nonempty subalgebra containing
\(\xi(X)\). Conversely every subalgebra containing
\(\xi(X)\) contains every term value, by induction
on terms, including constants.
Hence its image is precisely \(\mathbf S_{i,\xi}\).

Define \(h:\mathbf T(X)\to\mathbf S^*\) by
\[
 h(p)=(\widehat\xi(p))_{(i,\xi)}.
 \]
Each coordinate is in the specified generated
subalgebra. These are values of a fixed function,
not an unchosen family of elements.
A fixed term \(p_0\) gives \(h(p_0)\), proving
that this product has a point.
The index set is nonempty because one basis
algebra and one of its elements give a constant
valuation of the nonempty \(X\).
Coordinatewise operation clauses show \(h\)
is a homomorphism, and \(h(x)=z_x\).
Its image is exactly the subalgebra generated
by those \(z_x\): one inclusion follows by
closure of the image, the other by induction
on the terms whose values form the image.
Every restricted coordinate map has image
\(\widehat\xi(T)=S_{i,\xi}\), hence is onto that
generated subalgebra. This is the asserted
subdirectness, with no irreducibility conclusion.

Equality \(h(p)=h(q)\) says that every valuation
into every algebra of \(\mathfrak A\) equates
\(p,q\). Every homomorphism from the term algebra
is such an evaluation, by its universal property.
By (HSP.4) this is exactly
\((p,q)\in\theta_{\mathfrak A}(X)\).
Thus \([p]\mapsto h(p)\) is a well-defined
bijective homomorphism from the indicated quotient
to \(\mathbf F(\mathfrak A,X)\), and its
inverse preserves operations.
This proves (HSP.10), including the assertion
that the product realizes exactly the common laws.
Preservation also shows that all common laws
hold in its image under all valuations.

For clarity about the separate deductive bridge,
take \(U\) containing \(X\) and the variables of
the law basis \(B\). The deductive congruence
\(\theta_{B,U}\) is a congruence on
\(\mathbf T(U)\), and \(\mathbf T(X)\)
is its subalgebra. Its restriction \(R_{B,X}\)
is therefore a congruence on \(\mathbf T(X)\).
For \(p,q\in T(X)\), equational completeness gives
\[
 (p,q)\in R_{B,X}
 \quad\Longleftrightarrow\quad
 B\models_{\mathrm{alg}}p\approx q
 \quad\Longleftrightarrow\quad
 (p,q)\in\theta_{M(B)}(X).
 \]
In the middle equivalence, valuations on \(X\)
extend to \(U\) using one fixed element of the
nonempty carrier, and term values use only the
variables actually occurring. Satisfaction of
\(B\) is the global identity interpretation.
Thus the quotient of equated functions, under
the explicitly modern equality convention,
is the relative-free algebra for \(M(B)\).
The number of free generators is not a restriction
on the variable supports allowed in its law basis.

If \(B=\operatorname{Id}_{\mathfrak A}(W)\),
the renaming argument in
\cref{annex:p10-hsp-representation} equates the
relative-free kernels of \(M(B)\) and
\(\mathfrak A\) on arbitrary \(X\).
The all-valuations product and the equational
quotient therefore give isomorphic free algebras,
the finite-term version of original equation 10-1.
Taking \(X=\{x_a:a\in A\}\) for an algebra
obeying the common laws yields the onto map
\([p]\mapsto p^{\mathbf A}(x_a\mapsto a)\)
proved in the preceding annex. This establishes
the original converse construction at the
declared finitary scope.

For the family/law assertion, under the stated modern
equality bridge, let \(\Phi\) be a deductively closed
set of laws over the infinite alphabet \(W\).
Equational soundness and completeness give
\(\operatorname{Id}_{M(\Phi)}(W)=\Phi\):
soundness gives one inclusion, and any identity
true in every model is deducible from \(\Phi\),
hence already belongs to \(\Phi\).
For an H/S/P-closed nonempty family \(\mathfrak F\),
the HSP characterization gives
\(M(\operatorname{Id}_{\mathfrak F}(W))=\mathfrak F\).
These are exactly the two inverse assignments
in (HSP.11) with the declared source operand types.
Every resulting algebra family can in fact be
generated from a set basis consisting of one algebra.
Take \(\mathbf F_{\mathfrak F}(\bar W)\).
Its membership in the H/S/P-closed \(\mathfrak F\)
follows from the explicit witness embedding.
It has the same identities over \(W\) as
\(\mathfrak F\): one inclusion is preservation;
if a law holds in this free algebra, evaluating
it at the canonical generator classes puts its
term pair in \(\theta_{\mathfrak F}(W)\), and
(HSP.4) proves the law in every member of
\(\mathfrak F\). The HSP characterization therefore
gives \(\mathfrak F=HSP(\{\mathbf F_{\mathfrak F}(\bar W)\})\).
This supplies a set basis and avoids treating all
generator cardinalities as a set-indexed product.
Increasing an algebra family can only remove common
laws, and increasing a law family can only remove
models. The inverse assignments therefore reverse
inclusion in both directions.

The law families form a complete lattice: intersections
of deductively closed theories are closed, and the
join is deductive closure of their union. The top
is all identities and the bottom is the identities
deducible from no premises. This is a lattice on
a subset of the power set of \(T(W)^2\), so the
alphabet/set-size requirement is satisfied.
The inverse inclusion-reversing maps transport that
lattice to the algebra families, interchanging
meets with joins. This proves the asserted dual
isomorphism rather than assuming it from a name.

Finally suppose the basis has finite carrier sizes
\(a_i\) and \(X\) has \(m\) elements.
There are exactly \(a_i^m\) maps \(X\to A_i\).
Each corresponding \(\mathbf S_{i,\xi}\)
has at most \(a_i\) elements, so its entire product
has at most \(\prod_i a_i^{a_i^m}\) elements.
The generated diagonal subalgebra has no more
elements than this product.
An algebra generated by \(m\) elements and
obeying the common laws is an onto image of that
relative-free algebra, by the established generator
extension and target-quotient argument.
Its size is therefore bounded by (HSP.12).
No equality of the generated subalgebra with the
whole product is asserted.

It is not an assertion that finitary deductions
also prove the original's unrestricted
ordinal-arity/rank formulation.
\end{proof}
\end{annexitem}
